\subsection{在逆极限上的应用}

命题 5. 设 \((\mathbf{G}_{\alpha}, f_{\alpha \beta})\) 是豪斯多夫拓扑群的一个逆系统，使得诸 \(f_{\alpha \beta}\) 是核紧的满射严格态射。则对每个 \(\alpha \in I\) ，典范映射 \(f_{\alpha}\) （它把 \(\mathbf{G} = \varprojlim \mathbf{G}_{\alpha}\) 映入 \(\mathbf{G}_{\alpha}\) ）是一个核紧的满射严格态射。

\(f_{\alpha}\) 是满射的且其核是紧的这两件事是第一章 § 9 第 6 节命题 8 推论 1 的推论。剩下的只是证明 \(f_{\alpha}\) 是一个严格态射。用 \(e\) （相应地，\(e_{\alpha}\) ）表示 G （相应地，\(G_{\alpha}\) ）的单位元。G 中 \(e\) 的每个邻域 V 包含形如 \(\overline{f}_{\beta}^{-1}(V_{\beta})\) 的集合，这里 \(V_{\beta}\) 是 \(e_{\beta}\) 在 \(G_{\beta}\) 中的一个邻域，且我们可以假设 \(\beta \geqslant \alpha\) ；由于 \(f_{\alpha \beta}\) 是满射的严格态射，\(f_{\alpha \beta}(V_{\beta})\) 是 \(e_{\alpha}\) 在 \(G_{\alpha}\) 中的一个邻域，又由于 \(f_{\beta}\) 是满射的，我们有 \(V_{\beta} \subset f_{\beta}(V)\) ，由此

\[
f _ {\alpha} (\mathbf {V}) = f _ {\alpha \beta} (f _ {\beta} (\mathbf {V})) \supset f _ {\alpha \beta} (\mathbf {V} _ {\beta});
\]

这表明 \(f_{\alpha}(\mathbf{V})\) 是 \(e_{\alpha}\) 在 \(\mathbf{G}_{\alpha}\) 中的一个邻域。

若 \(\mathrm{H}_{\alpha}=\overline{f}_{\alpha}^{-1}(e_{\alpha})\) ，则每个 \(H_{\alpha}\) 是 G 的一个紧正规子群；诸 \(H_{\alpha}\) 显然满足第 3 节的条件 (AP)；且 \(G_{\alpha}\) 可与 \(G/H_{\alpha}\) 等同。特别地，第 3 节的命题 3 与命题 4 可应用于 G 与诸 \(H_{\alpha}\) 。

推论 1. 设 \((\mathbf{G}_{\alpha}, f_{\alpha \beta})\) 是满足命题 5 的假设的拓扑群逆系统；设 \((\mathbf{G}_{\alpha}', f_{\alpha \beta}')\) 是拓扑群的一个逆系统，且对每个 \(\alpha\) ，设 \(u_{\alpha}: \mathbf{G}_{\alpha} \to \mathbf{G}_{\alpha}'\) 是核紧的满射严格态射，使得诸 \(u_{\alpha}\) 构成映射的一个逆系统。则 \(u = \varprojlim u_{\alpha}\) 是 \(\mathbf{G} = \varprojlim \mathbf{G}_{\alpha}\) 到 \(\mathbf{G}' = \varprojlim \mathbf{G}_{\alpha}'\) 上的一个严格态射，且其核是紧的。

设 \(N_{\alpha}\) 是 \(u_{\alpha}\) 的核。\(L_{\alpha} = \overline{f}_{\alpha}^{-1}(N_{\alpha})\) 因而是满射严格态射 \(v_{\alpha} = u_{\alpha} \circ f_{\alpha}: G \to G'_{\alpha}\) 的核；由于 \(L_{\alpha}/H_{\alpha}\) 同构于 \(N_{\alpha}\) （§ 2，第 7 节，命题 20），\(L_{\alpha}\) 是 G 的一个紧正规子群（§ 4，第 1 节，命题 2 的推论 2）。u 的核 L 是诸 \(L_{\alpha}\) 的交。用 \(\varphi\) 表示典范映射 \(G \to G/L\) ；则我们可以写 \(v_{\alpha} = w_{\alpha} \circ \varphi\) ，其中 \(w_{\alpha}\) 是 G/L 到 \(G'_{\alpha}\) 上的一个严格态射，其核是 \(L_{\alpha}/L\) 。由于诸 \(L_{\alpha}/L\) 的交是 G/L 的单位元，且诸 \(L_{\alpha}/L\) 组成一个滤基并且是紧的，这个滤基收敛于 G/L 的单位元（第一章，§ 9，第 1 节，定理 1 的推论）。于是第 3 节的命题 2 表明 \(w = \varprojlim w_{\alpha}\) 是 G/L 到 \(G'\) 上的一个同构；由此推出 \(w \circ \varphi\) 是 G 到 \(G'\) 上的一个严格态射，其核为 L。但显然 \(u = w \circ \varphi\) ，这样推论就证明了。

推论 2. 设 \((\mathbf{G}_{\alpha}, f_{\alpha \beta})\) 是满足命题 5 的条件的拓扑群逆系统，并设 \(\mathbf{G}'\) 是一个拓扑群，其中存在一个邻域 \(\mathbf{V}'\) （单位元 \(e'\) 的），它不含 \(\mathbf{G}'\) 的除 \(\{e'\}\) 外的任何子群。那么，若 \(v: \mathbf{G} \to \mathbf{G}'\) 是任一连续同态，则存在指标 \(\alpha \in \mathbf{I}\) 与连续同态 \(v_{\alpha}: \mathbf{G}_{\alpha} \to \mathbf{G}'\) ，使得 \(v = v_{\alpha} \circ f_{\alpha}\) 。

因为由于 \(\overline{v}^{-1}(\mathbf{V}^{\prime})\) 是 \(e\) 在 \(\mathbf{G}\) 中的一个邻域，存在指标 \(\alpha\) 与一个邻域 \(\mathbf{V}_{\alpha}\) （\(e_{\alpha}\) 在 \(\mathbf{G}_{\alpha}\) 中的），使得 \(\overline{f}_{\alpha}^{-1}(\mathbf{V}_{\alpha}) \subset \overline{v}^{-1}(\mathbf{V}^{\prime})\) 。因而 \(v(\mathbf{H}_{\alpha}) \subset \mathbf{V}^{\prime}\) ，且由于 \(v(\mathbf{H}_{\alpha})\) 是 \(\mathbf{G}^{\prime}\) 的一个子群，由此推出 \(v(\mathbf{H}_{\alpha}) = \{e^{\prime}\}\) 。由于 \(f_{\alpha}\) 可与典范映射 \(\mathbf{G} \to \mathbf{G}/\mathbf{H}_{\alpha}\) 等同，本推论是连续同态的典范分解（§ 2，第 8 节）的推论。

